% Source: https://artofproblemsolving.com/community/c3h1973385_guys_help_its_moems
% Stable identity: aops-topic-1973385-q01
A number has an odd number of factors exactly when it is a perfect square, so statement 4 says $N+5$ is a perfect square. Statement 3 says $1<N<49$.

Statements 1 and 2 are both true only if $N$ and $N+3$ are both prime; one of them is even, so it must be $2$, which forces $N=2$. For $N=2$: statements 1, 2, 3 are true and $N+5=7$ is not a square, so exactly three are true. This is the only case with statements 1 and 2 both true.

Otherwise exactly one of statements 1, 2 is true and statements 3 and 4 are both true, so $1<N<49$ and $N+5\in\{9,16,25,36,49\}$, i.e. $N\in\{4,11,20,31,44\}$.
\begin{itemize}
\item $N=4$: $4$ is not prime, $7$ is prime. Exactly three true.
\item $N=11$: $11$ is prime, $14$ is not. Exactly three true.
\item $N=20$: $20$ is not prime, $23$ is prime. Exactly three true.
\item $N=31$: $31$ is prime, $34$ is not. Exactly three true.
\item $N=44$: $44$ is not prime, $47$ is prime. Exactly three true.
\end{itemize}
So $N\in\{2,4,11,20,31,44\}$, and the answer is $\boxed{6}$.
